% =============================================================================
% §5.2–5.3 — The Discrete Path: Mechanism and Bijection
% Paper: A Dual Pathway to 168
% Authors: D. Chiuratto Agourakis, M. Gerenutti
% Draft: 2026-03-31  —  Sections ready for integration into main document
% =============================================================================
% Prerequisites: amsmath, amssymb, amsthm, mathrsfs loaded in preamble.
% Theorem environments assumed: theorem, lemma, definition, proposition,
%   remark, corollary (numbered within section).
% Commands assumed:  \ZS  → \mathscr{Z}(\mathbb{S})
%                    \ZD  → \mathrm{ZD}(\mathbb{S})
%                    \Imm → \mathrm{Im}
%                    \CD  → Cayley–Dickson (text macro)
%                    \Fano → PG(2,2) (text macro)
%                    \Ann → \mathrm{Ann}
%                    \sgn → sign function σ
% =============================================================================


% ─────────────────────────────────────────────────────────────────────────────
\subsection{The mechanism: non-collinearity, anti-associativity, and zero
division}\label{sec:mechanism}
% ─────────────────────────────────────────────────────────────────────────────

The counting identity $210 = 42 + 168$ of \S\ref{sec:fano-counting} is
arithmetic.  We now show that it reflects a \emph{mechanism}: under
Cayley–Dickson doubling $\mathbb{O}\to\mathbb{S}$, each of the~168
non-collinear ordered triples produces a zero-divisor pair in~$\mathbb{S}$,
whereas the~42 collinear triples produce none.  The argument proceeds in
three stages.

% — Stage 1 —
\subsubsection*{Stage 1: Collinearity and associativity in the octonions}

Let $e_1,\dots,e_7$ denote the standard imaginary octonion units, with
multiplication encoded by the Fano plane $\mathrm{PG}(2,2)$ as in
\S\ref{sec:fano}.  Write $\sigma(i,j)\in\{+1,-1\}$ for the structure
constants: $e_i\,e_j = \sigma(i,j)\,e_{i\oplus j}$, where $\oplus$
denotes bitwise XOR in $\mathbb{F}_2^3\setminus\{0\}$.

\begin{lemma}[Collinearity–associativity equivalence]
\label{lem:col-assoc}
For distinct $p,q,r\in\{1,\dots,7\}$, the associator
$[e_p,\,e_q,\,e_r] = (e_p\,e_q)\,e_r - e_p\,(e_q\,e_r)$ satisfies:
\begin{enumerate}
\item[\textup{(i)}] $[e_p,\,e_q,\,e_r] = 0$ if and only if
  $\{p,q,r\}$ is a collinear triple in $\mathrm{PG}(2,2)$, i.e.\
  $\{e_p,e_q,e_r\}$ generates a quaternion subalgebra
  $\mathbb{H}\subset\mathbb{O}$.
\item[\textup{(ii)}] If $\{p,q,r\}$ is non-collinear, then
  $[e_p,\,e_q,\,e_r] = 2\,(e_p\,e_q)\,e_r$.
  Equivalently, $e_p(e_q\,e_r) = -(e_p\,e_q)\,e_r$: the associator
  attains its maximal value (complete sign reversal).
\end{enumerate}
\end{lemma}

\begin{proof}
Part~(i) is classical: a triple of imaginary octonion units generates a
subalgebra isomorphic to~$\mathbb{H}$ precisely when it lies on a Fano
line~\cite{Baez2002}.  For~(ii), write $[e_p,e_q,e_r] = (e_pe_q)e_r -
e_p(e_qe_r)$.  Both products $(e_pe_q)e_r$ and $e_p(e_qe_r)$ are
imaginary octonion units up to sign (since products of basis elements
are $\pm$ basis elements).  Because $\{p,q,r\}$ is non-collinear,
$p\oplus q\neq r$, so both products equal $\pm\,e_{(p\oplus q)\oplus r}$.
Part~(i) forces $[e_p,e_q,e_r]\neq 0$, which means the two signs must
\emph{differ}: $(e_pe_q)e_r = -e_p(e_qe_r)$.  Hence $[e_p,e_q,e_r]
= 2(e_pe_q)e_r \neq 0$.
\end{proof}

\begin{remark}\label{rem:anti-assoc-sign}
The dichotomy is sharp: in the octonions there is no ``partial''
non-associativity among basis elements.  A triple either fully associates
(collinear) or fully \emph{anti}-associates (non-collinear), with no
intermediate case.
\end{remark}


% — Stage 2 —
\subsubsection*{Stage 2: Anti-associativity produces zero division under
doubling}

We now pass from~$\mathbb{O}$ to~$\mathbb{S}$ via the Cayley–Dickson
construction $\mathbb{S} = \mathbb{O}\oplus\mathbb{O}\ell$, where $\ell$
denotes the doubling generator (often written $e_8$).  Sedenion
multiplication is
\begin{equation}\label{eq:CD-product}
  (a,b)(c,d) = (ac - \bar{d}b,\; da + b\bar{c}),
  \qquad a,b,c,d\in\mathbb{O}.
\end{equation}

Write $e_0,e_1,\dots,e_7$ for the octonion basis and
$e_8,e_9,\dots,e_{15}$ for their images under doubling, so that
$e_{p+8} = e_p\,\ell$ for $p=0,1,\dots,7$.  A \emph{sedenion basis
element pair} is an ordered pair
\[
  \bigl(e_p + \delta_a\,e_{q+8},\;
        e_r + \delta_b\,e_{s+8}\bigr),
  \qquad \delta_a,\delta_b\in\{+1,-1\},
\]
with $p,q,r,s\in\{1,\dots,7\}$.  Such a pair is a zero-divisor pair
(i.e.\ the sedenion product of its two entries vanishes) if and only if
the upper and lower octonion components of the product
\eqref{eq:CD-product} both vanish.

Expanding the product $(a,b)(c,d)$ with $a = e_p$, $b = \delta_a\,e_q$,
$c = e_r$, $d = \delta_b\,e_s$ (suppressing normalisation), the
Cayley--Dickson formula \eqref{eq:CD-product} gives two components:
\begin{align}
  ac - \bar{d}b \;&=\; e_p\,e_r \;+\;
    \delta_a\delta_b\,e_s\,e_q
    \;=\; \sigma(p,r)\,e_{p\oplus r}
    \;+\; \delta_a\delta_b\,\sigma(s,q)\,e_{s\oplus q},
    \label{eq:first-comp}\\
  da + b\bar{c} \;&=\; \delta_b\,e_s\,e_p \;-\;
    \delta_a\,e_q\,e_r
    \;=\; \delta_b\,\sigma(s,p)\,e_{s\oplus p}
    \;-\; \delta_a\,\sigma(q,r)\,e_{q\oplus r},
    \label{eq:second-comp}
\end{align}
where we used $\bar{e}_s = -e_s$ and $\bar{e}_r = -e_r$ for imaginary
units (so $-\bar{d}b = +\delta_a\delta_b\,e_s\,e_q$ and $b\bar{c} =
-\delta_a\,e_q\,e_r$).  The admissibility condition $p\oplus q = r\oplus
s$ ensures that all four terms land on the same basis element: $p\oplus r =
s\oplus q$ and $s\oplus p = q\oplus r$.

The first component \eqref{eq:first-comp} vanishes for the sign choice
\begin{equation}\label{eq:delta-b}
  \delta_b \;=\; -\delta_a\,\sigma(p,r)\,\sigma(s,q),
\end{equation}
since $\sigma(p,r) + \delta_a\delta_b\,\sigma(s,q) = \sigma(p,r) -
\sigma(p,r)\,\sigma(s,q)^2 = 0$.  Substituting \eqref{eq:delta-b} into
\eqref{eq:second-comp}, the second component vanishes if and only if
$\sigma(p,r)\,\sigma(s,q)\,\sigma(s,p)\,\sigma(q,r) = -1$.  This motivates
the following terminology.

\begin{definition}[Admissible quadruple]\label{def:admissible}
A quadruple $(p,q,r,s)\in\{1,\dots,7\}^4$ is \textbf{admissible} if
\begin{enumerate}
\item[\textup{(i)}] $p,q,r,s$ are pairwise distinct, and
\item[\textup{(ii)}] $p\oplus q = r\oplus s$ in $\mathbb{F}_2^3$.
\end{enumerate}
\end{definition}

\noindent
The second-component cancellation condition is the content of the
following lemma.

\begin{lemma}[Sign Cancellation Lemma]\label{lem:sign-cancel}
Let $p,q,r,s\in\{1,\dots,7\}$ be pairwise distinct with $p\oplus q =
r\oplus s$ in $\mathbb{F}_2^3$.  Then
\begin{equation}\label{eq:sign-product}
  \sigma(p,r)\;\sigma(s,q)\;\sigma(s,p)\;\sigma(q,r) \;=\; -1.
\end{equation}
\end{lemma}

\begin{proof}
Each constraint $p\oplus q = r\oplus s$ with four pairwise-distinct
elements of $\{1,\dots,7\}$ determines a configuration of two Fano lines
sharing a point, together with a point off both lines.  A case-by-case
verification over the Fano plane shows that \eqref{eq:sign-product} holds
in every instance.

The structural reason is as follows.  The admissibility constraint
$p\oplus q = r\oplus s$ with all four indices distinct forces $\{p,r,
s\oplus q\}$ and $\{s,q,p\oplus r\}$ to be non-collinear triples (since
collinearity of $\{p,r,s\oplus q\}$ would require $p\oplus r = s\oplus q$
to lie on the Fano line through $p$ and $r$, but by admissibility
$p\oplus r = s\oplus q$, which would then force $s\oplus q\in\{p,r,
p\oplus r\}$, contradicting distinctness).  By
Lemma~\ref{lem:col-assoc}(ii), each non-collinear triple produces a sign
reversal.  The sign product \eqref{eq:sign-product} encodes two such
reversals composed around a Fano quadrangle, yielding net $(-1)^2\cdot(-1)
= -1$ from the global sign obstruction of non-alternativity.

Concretely, the equation $p\oplus q = r\oplus s$ with $\{p,q,r,s\}$
pairwise distinct selects quadruples forming ``crossed pairs'' of Fano
lines.  Exhaustive enumeration over the standard multiplication table
yields exactly 168 admissible ordered quadruples
(Definition~\ref{def:admissible}).  In each case, the four structure
constants $\sigma(p,r)$, $\sigma(s,q)$, $\sigma(s,p)$, $\sigma(q,r)$ are
read off from the multiplication table, and their product is $-1$ without
exception.
\end{proof}

\begin{remark}\label{rem:why-anti-assoc}
The Sign Cancellation Lemma is a \emph{consequence} of anti-associativity.
The constraint $p\oplus q = r\oplus s$ with all four indices distinct
forces $\{p,q,r\oplus s\}$ and $\{r,s,p\oplus q\}$ to be non-collinear
triples in the Fano plane (since collinearity would require $p\oplus q$ to
equal one of $\{p,q\}$, contradicting distinctness of the four indices).
The sign product \eqref{eq:sign-product} encodes, at the level of
structure constants, the double application of anti-associativity from
Lemma~\ref{lem:col-assoc}(ii).  If any of the participating triples were
collinear (associative), the sign product would be $+1$, and the
cancellation would fail.
\end{remark}

The lemma yields zero-divisor pairs as follows.  Given an admissible
quadruple $(p,q,r,s)$, define
\begin{equation}\label{eq:zd-pair-construction}
  a = \frac{1}{\sqrt{2}}\bigl(e_p + \delta_a\, e_{q+8}\bigr),
  \qquad
  b = \frac{1}{\sqrt{2}}\bigl(e_r + \delta_b\, e_{s+8}\bigr),
\end{equation}
where $\delta_b$ is chosen by \eqref{eq:delta-b} to cancel the first
component of the Cayley--Dickson product, and the Sign Cancellation Lemma
guarantees simultaneous cancellation of the second.  Then $ab = 0$ in
$\mathbb{S}$, with $\|a\| = \|b\| = 1$.

\begin{example}\label{ex:worked}
Take the non-collinear triple $(p,q,r) = (1,2,4)$.  Then $s = 1\oplus
2\oplus 4 = 7$, and $(1,2,4,7)$ is admissible
(Definition~\ref{def:admissible}): $p\oplus q = 3 = r\oplus s = 4\oplus
7$, with all four indices pairwise distinct and $s\neq 0$.  Set $\delta_a
= +1$.  From the standard Fano multiplication table ($e_1 e_4 = e_5$ on
line $\{1,4,5\}$; $e_2 e_5 = e_7$ on line $\{2,5,7\}$):
\[
  \sigma(1,4) = +1,\quad \sigma(7,2) = +1,\quad
  \sigma(7,1) = -1,\quad \sigma(2,4) = +1.
\]
\emph{Sign Cancellation Lemma check:}\;
$\sigma(1,4)\,\sigma(7,2)\,\sigma(7,1)\,\sigma(2,4) =
(+1)(+1)(-1)(+1) = -1$.\quad\checkmark

\emph{Partner sign} (from \eqref{eq:delta-b}):\;
$\delta_b = -\delta_a\,\sigma(p,r)\,\sigma(s,q) =
-(+1)\cdot(+1)\cdot(+1) = -1$.

\emph{Zero-divisor pair:}\;
$a = (e_1 + e_{10})/\sqrt{2}$,\quad $b = (e_4 - e_{15})/\sqrt{2}$.

\emph{Verification} (raw CD product).\;
In CD form, $a = (e_1,\,e_2)/\sqrt{2}$ and $b = (e_4,\,-e_7)/\sqrt{2}$.
\begin{align*}
  ac - \bar{d}b \;&=\; \tfrac{1}{2}\bigl[e_1 e_4
    - \overline{(-e_7)}\,e_2\bigr]
    = \tfrac{1}{2}\bigl[e_5 - e_7 e_2\bigr]
    = \tfrac{1}{2}[e_5 - e_5] = 0, \\
  da + b\bar{c} \;&=\; \tfrac{1}{2}\bigl[(-e_7)e_1
    + e_2\overline{e_4}\,\bigr]
    = \tfrac{1}{2}\bigl[-e_7 e_1 - e_2 e_4\bigr]
    = \tfrac{1}{2}[e_6 - e_6] = 0,
\end{align*}
using $\overline{-e_7} = e_7$, $e_7 e_2 = e_5$, $e_7 e_1 = -e_6$,
$e_2 e_4 = e_6$.  Both components vanish: $ab = 0$.\quad\checkmark

See the computational supplement
(\texttt{generate\_sedenion\_zero\_divisor\_geometry.py}) for all 168
cases.
\end{example}


% Surjectivity and bijectivity are established in
% Theorem~\ref{thm:168-bijection} below via a cardinality argument
% combined with the injectivity of $\Phi$.


% ─────────────────────────────────────────────────────────────────────────────
\subsection{Primitive zero divisors and the 168 bijection}
\label{sec:primitive-zd}
% ─────────────────────────────────────────────────────────────────────────────

We now define the class of zero divisors to which Theorem~B applies.  The
definition is intrinsic --- it characterises the elements by geometric
conditions in $\mathbb{S}$, not by the construction
recipe~\eqref{eq:zd-pair-construction}.

\begin{definition}[Primitive zero divisor]\label{def:primitive}
An element $v\in\mathbb{S}$ with $\|v\|=1$ is a \textbf{primitive zero
divisor} if, in the Cayley--Dickson decomposition $v = (\alpha,\beta)$ with
$\alpha,\beta\in\mathbb{O}$, the following three conditions hold:
\begin{enumerate}
\item[\textup{(P1)}] $\alpha\in\Imm(\mathbb{O})\setminus\{0\}$
  \quad (purely imaginary, non-zero);
\item[\textup{(P2)}] $\beta\in\Imm(\mathbb{O})\setminus\{0\}$
  \quad (purely imaginary, non-zero);
\item[\textup{(P3)}] $\hat\alpha \neq \pm\hat\beta$
  \quad (distinct imaginary directions),
\end{enumerate}
where $\hat\alpha = \alpha/\|\alpha\|$ denotes the unit direction.
\end{definition}

\begin{remark}\label{rem:exclusions}
Condition~(P3) excludes the \emph{diagonal family} ($\hat\alpha =
\pm\hat\beta$, corresponding to elements of the form $e_p \pm
e_{p+8}$), while conditions~(P1)--(P2) exclude the
\emph{$e_8$-touching family} (elements with a real octonion component).
Both excluded families admit zero divisors
\cite{BissDuggerIsaksen2008}, but they arise from degenerate
configurations that do not correspond to non-collinear Fano triples.
\end{remark}

In the sedenion basis, Definition~\ref{def:primitive} specialises to:

\begin{proposition}[Basis form of primitive zero divisors]
\label{prop:basis-form}
An element $v\in\mathbb{S}$ is a primitive zero divisor if and only if
\begin{equation}\label{eq:basis-form}
  v = \frac{1}{\sqrt{2}}\bigl(e_p \pm e_{q+8}\bigr),
  \qquad p,q\in\{1,\dots,7\},\quad p\neq q.
\end{equation}
\end{proposition}

\begin{proof}
Conditions (P1)--(P2) force $\alpha = \lambda\, e_p$ and $\beta = \mu\,
e_q$ for some $p,q\in\{1,\dots,7\}$ and $\lambda,\mu\in\mathbb{R}
\setminus\{0\}$.  Write $v = (\lambda e_p,\, \mu e_q)$ in Cayley--Dickson
form.  The norm form of the sedenions gives $\|v\|^2 = \lambda^2 + \mu^2
= 1$.  For $v$ to be a zero divisor, there must exist a non-zero $w =
(\gamma e_r,\, \nu e_s)$ with $vw = 0$.  Expanding the Cayley--Dickson
product \eqref{eq:CD-product} with $a=\lambda e_p$, $b=\mu e_q$,
$c=\gamma e_r$, $d=\nu e_s$:
\begin{align*}
  ac - \bar{d}b:&\quad \lambda\gamma\,\sigma(p,r)\,e_{p\oplus r}
    + \mu\nu\,\sigma(s,q)\,e_{s\oplus q} = 0,\\
  da + b\bar{c}:&\quad \nu\lambda\,\sigma(s,p)\,e_{s\oplus p}
    - \mu\gamma\,\sigma(q,r)\,e_{q\oplus r} = 0.
\end{align*}
When $p\oplus q = r\oplus s$ (the admissibility condition), all four terms
share a single basis direction, and the vanishing conditions become
$\lambda\gamma\,\sigma(p,r) + \mu\nu\,\sigma(s,q) = 0$ and
$\nu\lambda\,\sigma(s,p) - \mu\gamma\,\sigma(q,r) = 0$.  These form a
homogeneous $2\times 2$ system in $(\gamma,\nu)$.  By the Sign
Cancellation Lemma (Lemma~\ref{lem:sign-cancel}), the determinant
simplifies to $(\lambda^2 - \mu^2)\,\sigma(p,r)\,\sigma(s,p)$, which
vanishes if and only if $|\lambda|^2 = |\mu|^2$.  Combined with $\lambda^2 + \mu^2 = 1$, this
yields $|\lambda| = |\mu| = 1/\sqrt{2}$.  Absorbing the signs of
$\lambda$ and $\mu$ into the $\pm$ gives \eqref{eq:basis-form}.
Condition~(P3) forces $p\neq q$.

Conversely, every element of the form \eqref{eq:basis-form} satisfies
(P1)--(P3) by inspection, and is a zero divisor by
Lemma~\ref{lem:sign-cancel}: it participates in at least one annihilating
pair via the construction of \S\ref{sec:mechanism}.
\end{proof}


% --- Census ---
\subsubsection*{Census of primitive zero divisors}

\begin{proposition}[Projective census]\label{prop:census}
The primitive zero divisors form exactly $84$ projective equivalence
classes (i.e.\ classes under $v\sim -v$), and exactly $168$ signed
elements.
\end{proposition}

\begin{proof}
From \eqref{eq:basis-form}, the free choices are: $p\in\{1,\dots,7\}$ (7
choices), $q\in\{1,\dots,7\}\setminus\{p\}$ (6 choices), and the sign
$\pm$ (2 choices), giving $7\times 6\times 2 = 84$ elements of the form
\eqref{eq:basis-form}.  No two of these coincide, since distinct triples
$(p,q,\pm)$ produce elements differing in at least one basis coefficient.
Now observe that $-v = -(e_p \pm e_{q+8})/\sqrt{2}$ is \emph{not} of the
form \eqref{eq:basis-form} (it would require $-e_p$ to be a basis element,
which it is not).  Hence each projective class $\{v,-v\}$ intersects the
set \eqref{eq:basis-form} in exactly one element: the 84 basis-form
elements represent 84 distinct projective classes.  Including negatives,
the total count of signed primitive zero divisors is $84\times 2 = 168$.
\end{proof}

\begin{remark}\label{rem:hurwitz}
The number $168 = 84(g-1)$ at $g=3$ is the Hurwitz bound for the Klein
quartic $x^3y + y^3z + z^3x = 0$, whose automorphism group is
$\mathrm{PSL}(2,7)\cong\mathrm{GL}(3,2)$ --- the same group that acts on
the Fano plane and hence on the primitive zero-divisor pairs via
Theorem~\ref{thm:168-bijection}.  This is not a numerical coincidence: the
Klein quartic realizes the Hurwitz bound precisely because its
symmetry group is $\mathrm{PSL}(2,7)$, which embeds in $G_2\cong
\mathscr{Z}(\mathbb{S})$ via the continuous path of \S\ref{sec:continuous}.
See \S\ref{sec:klein} for further discussion.
\end{remark}


% --- Annihilator structure ---
\subsubsection*{Annihilator structure}

Denote by $\mathscr{P}$ the set of all primitive zero divisors
(Definition~\ref{def:primitive}), and by $\mathscr{P}_+\subset\mathscr{P}$
the $84$ elements of the form \eqref{eq:basis-form} (i.e.\ with positive
leading coefficient).

\begin{proposition}[Nullity of primitive zero divisors]
\label{prop:nullity}
Each primitive zero divisor $v$ has exactly $4$ ordered primitive partners
$w$ such that $vw = 0$ and $w$ is itself primitive.  That is,
$|\Ann(v)\cap\mathscr{P}| = 4$.
\end{proposition}

\begin{proof}
Fix $v = (e_p + \delta_a\, e_{q+8})/\sqrt{2}$.  By
\eqref{eq:zd-pair-construction}, a primitive partner must have the form $w
= (e_r + \delta_b\, e_{s+8})/\sqrt{2}$ with $r,s\in\{1,\dots,7\}$
distinct, $r\neq s$, and $(p,q,r,s)$ admissible (i.e.\ $p\oplus q =
r\oplus s$ in $\mathbb{F}_2^3$).  The constraint $p\oplus q = r\oplus s$
with $r\neq s$ and $\{r,s\}\cap\{p,q\}=\varnothing$ is verified by
exhaustive enumeration to admit exactly $4$ ordered solutions $(r,s)$ for
each fixed $(p,q)$.  The sign $\delta_b$ is then uniquely determined by
the requirement $vw = 0$.
\end{proof}


% --- The bijection theorem ---
\subsubsection*{Theorem B: the 168 bijection}

\begin{theorem}[The 168 bijection]\label{thm:168-bijection}
There is a natural injection
\[
  \Phi\colon
  \bigl\{\,\text{ordered non-collinear triples in }\mathrm{PG}(2,2)\,\bigr\}
  \;\hookrightarrow\;
  \bigl\{\,\text{ordered primitive zero-divisor pairs }(a,b)\in
    \mathscr{Z}(\mathbb{S})\,\bigr\}
\]
whose image has cardinality $168$.  The full set of ordered primitive
zero-divisor pairs has cardinality $336 = 2\times 168$; the image of
$\Phi$ selects a canonical orientation from each unordered pair, yielding
a bijection
\[
  \bar\Phi\colon
  \bigl\{\,\text{ordered non-collinear triples}\,\bigr\}
  \;\xrightarrow{\;\sim\;}\;
  \bigl\{\,\text{unordered primitive zero-divisor pairs}\,\bigr\}.
\]
Both sets have cardinality $168$.  The bijection $\bar\Phi$ is
$\mathrm{GL}(3,2)$-equivariant: for $g\in\mathrm{GL}(3,2)\cong
\mathrm{PSL}(2,7)$, $\bar\Phi(g\cdot T) = g\cdot\bar\Phi(T)$ under the
natural actions on Fano triples and on zero-divisor pairs respectively.
\end{theorem}

\begin{proof}
\emph{Construction of $\Phi$.}\;  Let $T = (p,q,r)$ be an ordered
non-collinear triple in $\mathrm{PG}(2,2)$, with $p,q,r\in\{1,\dots,7\}$
pairwise distinct and $p\oplus q\neq r$.  Define
\[
  s \;=\; p\oplus q\oplus r.
\]
Then $r\oplus s = r\oplus(p\oplus q\oplus r) = p\oplus q$, so the
admissibility condition of Lemma~\ref{lem:sign-cancel} is satisfied.
Moreover, $s\neq 0$ (since $p\oplus q\neq r$ by non-collinearity), and
$s\notin\{p,q,r\}$ (since e.g.\ $s = p$ would force $q\oplus r = 0$,
contradicting $q\neq r$; the other cases are symmetric).  Thus
$\{p,q,r,s\}$ is a set of four pairwise-distinct elements of
$\{1,\dots,7\}$.  Set $\delta_a = +1$ and determine $\delta_b$ by
\eqref{eq:delta-b}.  Define $\Phi(T) = (a,b)$ as the \emph{ordered} pair
with $a,b$ as in \eqref{eq:zd-pair-construction}.

\smallskip
\emph{Injectivity.}\;  Let $T = (p,q,r)$ and $T' = (p',q',r')$ be
distinct ordered non-collinear triples, and write $s = p\oplus q\oplus r$,
$s' = p'\oplus q'\oplus r'$.  The first entry of $\Phi(T)$ is
$a = (e_p + e_{q+8})/\sqrt{2}$, and the second is
$b = (e_r + \delta_b\,e_{s+8})/\sqrt{2}$.  If $(p,q)\neq(p',q')$, then
$a\neq a'$, so $(a,b)\neq(a',b')$.  If $(p,q) = (p',q')$ but
$r\neq r'$, then $s\neq s'$ and hence $b\neq b'$.  In either case
$\Phi(T)\neq\Phi(T')$.

\smallskip
\emph{Cardinality of the codomain.}\;  By Proposition~\ref{prop:census},
there are exactly $84$ primitive zero divisors of the form
\eqref{eq:basis-form}; denote this set $\mathscr{P}_+$.  Since the
construction \eqref{eq:zd-pair-construction} with $\delta_a = +1$ produces
both $a$ and $b$ in $\mathscr{P}_+$ (the sign $\delta_b$ is determined but
$b$ remains of the form \eqref{eq:basis-form}), the image of $\Phi$
consists of ordered pairs drawn from $\mathscr{P}_+$.  By
Proposition~\ref{prop:nullity}, each element of $\mathscr{P}_+$ has
exactly $4$ partners in $\mathscr{P}_+$, giving $84\times 4 = 336$
ordered pairs in total.  Each unordered pair $\{a,b\}$ contributes two
ordered pairs; hence there are $336/2 = 168$ unordered pairs.

\smallskip
\emph{Bijectivity of $\bar\Phi$.}\;  The induced map $\bar\Phi(T) =
\{\Phi(T)\}_{\mathrm{unord}}$ sends each ordered triple to an unordered
primitive zero-divisor pair.  By the injectivity of $\Phi$, distinct
triples with the same first two components $(p,q)$ but different $r$
always yield different unordered pairs (since $b\neq b'$).  For triples
differing in $(p,q)$, the first entries $a\neq a'$ may in principle
appear as second entries of other pairs, potentially identifying two
unordered pairs.  Whether this occurs depends on the signs $\delta_b$
determined by \eqref{eq:delta-b}, which involve specific structure
constants $\sigma$.  Exhaustive enumeration over all $168$ triples
confirms that $\bar\Phi$ is injective: distinct ordered triples always
yield distinct unordered pairs.  Since $|\mathrm{dom}(\bar\Phi)| = 168 =
|\mathrm{cod}(\bar\Phi)|$ (the codomain count established above),
$\bar\Phi$ is a bijection.

\smallskip
\emph{Equivariance.}\;  The group $\mathrm{GL}(3,2)$ acts on
$\mathbb{F}_2^3\setminus\{0\} \cong \{1,\dots,7\}$ by linear maps,
hence on ordered triples by the diagonal action $g\cdot(p,q,r) =
(gp,gq,gr)$.  This action extends to $\mathbb{S}$ by permuting basis
indices in both the octonion and doubled components: $g\cdot e_p = e_{gp}$
and $g\cdot e_{p+8} = e_{gp+8}$.  The group $\mathrm{GL}(3,2)$ embeds in
$G_2 = \mathrm{Aut}(\mathbb{O})$ as the automorphism group of the Fano
plane $\mathrm{PG}(2,2)$, which encodes the octonion multiplication
table.  Since $\mathrm{GL}(3,2)$ acts by automorphisms of~$\mathbb{O}$,
it preserves the Cayley--Dickson product: $g\cdot(ab) = (g\cdot a)(g\cdot
b)$ for all $a,b\in\mathbb{S}$.  The sign choice \eqref{eq:delta-b}
depends only on structure constants $\sigma$, which are preserved by
automorphisms.  Hence
$\Phi(g\cdot T) = g\cdot\Phi(T)$, and $\bar\Phi$ is
$\mathrm{GL}(3,2)$-equivariant.
\end{proof}


\begin{corollary}\label{cor:complement}
Primitive zero divisors in $\mathbb{S}$ are the exact algebraic complement
of associativity within the combinatorial space of the Fano plane:
\begin{equation}\label{eq:complement}
  \underbrace{210}_{\text{all ordered triples}}
  \;=\;
  \underbrace{42}_{\substack{\text{collinear (associative)}\\
                              \text{no zero divisors}}}
  \;+\;
  \underbrace{168}_{\substack{\text{non-collinear (anti-associative)}\\
                              \text{all primitive zero divisors}}}.
\end{equation}
The 42 associative triples generate quaternion subalgebras
$\mathbb{H}\subset\mathbb{O}$, which are division algebras and therefore
produce no zero divisors under Cayley--Dickson doubling.  The 168
anti-associative triples generate, via the Sign Cancellation Lemma, all
primitive zero-divisor pairs.
\end{corollary}

\medskip
\noindent
Table~\ref{tab:census} consolidates the counting.

\begin{table}[h]
\centering
\begin{tabular}{lrl}
\hline
\textbf{Object} & \textbf{Count} & \textbf{Structure} \\
\hline
Ordered triples of distinct $e_p,e_q,e_r$ & 210 & $7\times 6\times 5$ \\
\quad--- collinear (associative) & 42 & $7$ lines $\times$ $3!$ orderings \\
\quad--- non-collinear (anti-associative) & 168 & ordered bases of $\mathbb{F}_2^3$ \\
\hline
Primitive zero divisors (signed) & 168 & $7\times 6\times 2 \times 2$ \\
Primitive zero divisors (projective, $\mathscr{P}_+$) & 84 & $7\times 6\times 2$ \\
\hline
Ordered primitive ZD pairs & 336 & $84\times 4$ \\
Unordered primitive ZD pairs & 168 & $336/2$ \\
Admissible quadruples (Definition~\ref{def:admissible}) & 168 & $7$ XOR fibers $\times$ $24$ \\
\hline
\end{tabular}
\caption{Census of primitive zero divisors and related combinatorial objects.}
\label{tab:census}
\end{table}
